\documentclass[10pt]{beamer}

%\usepackage[backend=bibtex,firstinits=true,style=verbose-inote,citestyle=authortitle]{biblatex}
\usepackage{bm}
\usepackage{graphicx}
\usepackage{subcaption}
\usepackage{amsmath}
\usepackage{amsfonts}
\usepackage{makecell}
\usepackage{filecontents}
\usepackage{biblatex}
% \newcommand{\expect}[2][]{
\ifthenelse{\equal{#1}{}}{
\mathbb{E}\left[#2\right]
}{
\underset{#1}{\mathbb{E}}\left[#2\right]
}}

\newcommand{\cov}[2][]{
\ifthenelse{\equal{#1}{}}{
\text{Cov}\left[#2\right]
}{
\underset{#1}{\text{Cov}}\left[#2\right]
}}


\newcommand{\var}[2][]{
\ifthenelse{\equal{#1}{}}{
\text{Var}[#2]
}{
\underset{#1}{\text{Var}}[#2]
}}

\newcommand{\loss}[2][]{
\ifthenelse{\equal{#1}{}}{
\mathcal{L}(#2)
}{
\mathcal{L}_{#1}(#2)
}}

\newcommand{\kl}[2]{
\text{D}_\text{KL}[#1 \parallel #2]
}

\newcommand{\R}{\mathbb{R}}
%\newcommand{\Prob}{\mathbb{P}}

\newcommand{\1}[1]{\mathds{1}\{#1\}}


%\usecolortheme{dolphin}
\setbeamertemplate{navigation symbols}{}
\setbeamertemplate{section in toc}{\inserttocsectionnumber.~\inserttocsection}

\begin{filecontents*}{references.bib}
@misc{FixMatch,
Author = {Kihyuk Sohn and David Berthelot and Chun-Liang Li and Zizhao Zhang and Nicholas Carlini and Ekin D. Cubuk and Alex Kurakin and Han Zhang and Colin Raffel},
Title = {FixMatch: Simplifying Semi-Supervised Learning with Consistency and Confidence},
Year = {2020},
Eprint = {arXiv:2001.07685},
}
\end{filecontents*}

\addbibresource{references.bib}


\title{FixMatch: Simplifying Semi-Supervised Learning with
Consistency and Confidence\footnote{\citepaper{FixMatch}}}
%\subtitle{}
%\author{Ivan Skorokhodov}
%\date{}
%\logo{\includegraphics[height=1cm]{images/ipavlov-logo.png}}

\newcommand{\citepaper}[1]{\citetitle{#1} by \citeauthor{#1}}

%\graphicspath{{./images}}

%\usetheme{lucid}
\begin{document}

\begin{frame}
    \titlepage
\end{frame}

\begin{frame}{Overview}
    FixMatch:
    \begin{itemize}
        \item\pause is a strong Self-Supervised Learning approach
        \item\pause has a very simple loss function which has an interesting interpretation
        \item\pause has quite complex augmentation strategy
        \item\pause achieves 88.61\% accuracy on CIFAR-10 with 4 labels per class!
    \end{itemize}
\end{frame}

\begin{frame}{Notation}
    \begin{itemize}
        \item\pause Imagine we have two datasets of images:
            \begin{itemize}
                \item\pause Labeled images $X^l = \left\{\left(x_{n}^l, p_{n}\right)\right\}_{n=1}^N$
                \item\pause Unlabeled images $X^u = \left\{x_{n}^u \right\}_{n=1}^{k \times N}$, i.e. $X^u$ is $k$ times larger than $X^l$
                \item\pause $p_n$ is a one-hot class label
            \end{itemize}
        \item\pause Let $H(p, q)$ be a cross-entropy between $p$ and $q$.
        \item\pause We also have augmentation functions:
        \begin{itemize}
            \item\pause $\alpha(x)$ is a weak (i.e. simple) augmentation: random horizontal flipping and translations
            \item\pause $\mathcal{A}(x)$ is a strong (i.e. sophisticated) augmentation: color inversion, translation, contrast adjustment, etc
        \end{itemize}
    \end{itemize}
\end{frame}

\begin{frame}{FixMatch}
    \pause
    FixMatch loss is a combination of two losses:
    \begin{equation}
        \mathcal{L}_\text{FM} = \mathcal{L}_\text{cls} + \lambda_\text{pl}\mathcal{L}_\text{pl}
    \end{equation}
    \begin{itemize}
        \item\pause $\mathcal{L}_\text{cls}$ is a usual cross-entropy classification loss on $X^l$ dataset.
        \item\pause $\mathcal{L}_\text{pl}$ is a \textit{pseudo-labeling loss}, i.e. a cross-entropy loss that uses synthetic targets produced by our model
    \end{itemize}
\end{frame}

\begin{frame}{Pseudo-labeling loss}
\pause
Pseudo-labeling loss $\mathcal{L}_\text{pl}$ equals:

\begin{equation}
\mathcal{L}_\text{pl} = \frac{1}{kN} \sum_{n=1}^{kN} 1\left[\max \left(\bar{q}_n\right) \geq \tau\right] H\left(\bar{q}_{n}', p_m\left(y | \mathcal{A}\left(x_{n}\right)\right)\right)
\end{equation}

Where:
\begin{itemize}
    \item $\bar{q}_n = p_m(y | \alpha(x_n))$, i.e. class probabilities for weakly augmented $x_n$
    \item $\bar{q}_n' = \arg\max \bar{q}_n$, i.e. a pseudo label (one-hot)
    \item $\tau$ is a hyperparameter
\end{itemize}

\pause
Algorithm:
\begin{itemize}
    \item\pause Pick an unlabeled image $x_n$, produce two augmentated versions: $\bar{x}_n = \alpha(x)_n$ and $\tilde{x}_n = \mathcal{A}(x_n)$
    \item\pause Compute class probabilities $\bar{q}_n$ and $\tilde{q}_n$ for $\bar{x}_n$ and $\tilde{x}_n$
    \item\pause Pick only examples with confident class probabilities for weakly-augmented images
    \item\pause Cross-entropy term forces the model to give the same predictions for a weakly-augmented and a strongly-augmented images
\end{itemize}
\end{frame}

\begin{frame}{An interesting property of this pseudo-labeling loss}
    \begin{itemize}
        \item\pause Previously used variants of pseudo-labelling loss required tuning of $\lambda_\text{pl}$ weight during training and gradually increase it.
        \item\pause In FixMatch model becomes gradually more confident in new images and authors omit $\lambda_\text{pl}$ whatsoever!
        \item\pause So we get a curriculum learning out of the box!
    \end{itemize}
\end{frame}

\begin{frame}{Notes}
\begin{itemize}
    \item\pause FixMatch provides very strong scores
    \item\pause Has a very simple loss with an interesting side-effect of curriculum learning
    \item\pause A disadvantage: strong augmentations are based on CutOut, CTAugment, etc and seem \textit{very} sophisticated
    \item\pause Strongly beats SotA in many setups:
\end{itemize}

\begin{figure}
    \centering
    \includegraphics[width=\textwidth]{images/fixmatch-results.png}
    %\caption{FixMatch results}
    %\label{fig:results}
\end{figure}
\end{frame}

\end{document}
